% 原书物理页=印刷页20--23；2.2 全节（起于 p20 中部的 2.2 标题；p20 上部的内容属 2.1，已见 ch02-s01.tex）.
\section{Lie 代数的泛包络代数}\label{sec:ssl-2-2}
同本书其他地方一样，这里所有结合 $K$-代数都假定含单位元. 回忆一下，这样的结合代数可以按括号 $[a, b] = ab - ba$ 视为一个 $K$-Lie 代数.
\begin{theorem}\label{thm:ssl-2-2-1}
设 $\mathfrak g$ 是一个 $K$-Lie 代数. 存在一个由
\begin{itemize}[label={}]
\item 一个结合 $K$-代数 $\mathcal{U}\mathfrak g$，以及
\item 一个 Lie 代数态射 $\mathfrak g \longrightarrow \mathcal{U}\mathfrak g$
\end{itemize}
组成的二元组 $(\mathcal{U}\mathfrak g, \iota)$，它满足下面的\newterm{ssl-universal-property}{泛性质}{universal property}：对每个结合 $K$-代数 $A$，映射
\[\operatorname{Hom}_{K\text{-alg}}(\mathcal{U}\mathfrak g, A) \longrightarrow \operatorname{Hom}_{K\text{-Lie}}(\mathfrak g, A),\qquad f \longmapsto f \circ \iota\]
是双射.
\end{theorem}
泛性质保证了二元组 $(\mathcal{U}\mathfrak g, \iota)$ 在唯一同构意义下的唯一性. 对于喜欢范畴论的读者，可以把这个定理转述为下面两条中的任意一条：
\begin{itemize}[label={}]
\item 对每个 $\mathfrak g$，函子 $A \mapsto \operatorname{Hom}_{K\text{-Lie}}(\mathfrak g, A)$ 是可表的；
\item 把每个结合代数映到其底层 Lie 代数的函子具有左伴随.
\end{itemize}
后一种表述看起来更强，因为它包含了 $\mathfrak g \mapsto \mathcal{U}\mathfrak g$ 的函子性，但两者实际上是等价的.
\begin{proof}
考虑 $T\mathfrak g = \bigoplus_{n \in \mathbb{N}} T^{n}\mathfrak g$，其中 $T^{n}\mathfrak g = \mathfrak g^{\otimes n} = \mathfrak g \otimes_{K} \mathfrak g \otimes_{K} \cdots \otimes_{K} \mathfrak g$. 我们在这个 $K$-向量空间上赋予拼接积：它把 $T^{n}\mathfrak g \times T^{m}\mathfrak g$ 送到 $T^{n+m}\mathfrak g$，并且在初等张量上由
\[(x_{1} \otimes \cdots \otimes x_{n})(y_{1} \otimes \cdots \otimes y_{m}) = x_{1} \otimes \cdots \otimes x_{n} \otimes y_{1} \otimes \cdots \otimes y_{m}\]
给出. 这样就得到一个\newterm{ssl-graded}{分次}{graded}的结合 $K$-代数 $T\mathfrak g$，称为 $\mathfrak g$ 的\newterm{ssl-tensor-algebra}{张量代数}{tensor algebra}.
注意这一构造其实只依赖于 $\mathfrak g$ 底层的 $K$-向量空间，因此适用于任何 $K$-向量空间. 代数 $T\mathfrak g$ 配上映射 $\mathfrak g = T^{1}\mathfrak g \hookrightarrow T\mathfrak g$ 后满足下面的泛性质：对每个结合 $K$-代数 $A$，映射
\[\operatorname{Hom}_{K\text{-alg}}(T\mathfrak g, A) \longrightarrow \operatorname{Hom}_{K}(\mathfrak g, A)\]
是双射. 换句话说，函子 $V \mapsto TV$ 是``把代数映到其底层向量空间''的函子的左伴随. 这一点的验证很容易，但我们不给出细节，可参看任何一本代数书（例如 Lang 的）.
现在考虑由形如
\[x \otimes y - y \otimes x - [x, y] \qquad \in T^{\leqslant 2}\mathfrak g\]
的元素生成的双边理想 $I$. 注意它不是齐次理想. 最后置 $\mathcal{U}\mathfrak g := T\mathfrak g / I$. 这是一个结合代数，不再是分次的，但仍然是\newterm{ssl-filtered}{滤过}{filtered}的：若 $\mathcal{U}\mathfrak g^{\leqslant n}$ 是 $T^{\leqslant n}\mathfrak g$ 的像，则 $\mathcal{U}\mathfrak g^{\leqslant n}\,\mathcal{U}\mathfrak g^{\leqslant m} \subset \mathcal{U}\mathfrak g^{\leqslant n + m}$. 按构造，有
\[\operatorname{Hom}_{K\text{-alg}}(\mathcal{U}\mathfrak g, A) = \bigl\{ f \in \operatorname{Hom}_{K\text{-alg}}(T\mathfrak g, A),\ \forall x, y \in \mathfrak g,\ f(x)f(y) - f(y)f(x) = f([x, y]) \bigr\}.\]
于是，配上 $\iota : \mathfrak g \hookrightarrow T\mathfrak g \twoheadrightarrow \mathcal{U}\mathfrak g$ 的 $\mathcal{U}\mathfrak g$ 的所述泛性质可由 $T\mathfrak g$ 的泛性质导出.
\end{proof}
\textbf{应用.}\quad 若 $(V, \rho)$ 是 $\mathfrak g$ 的一个 $K$-表示，则态射 $\rho : \mathfrak g \longrightarrow \mathfrak{gl}(V)$ 来自（唯一的）代数态射 $\mathcal{U}\mathfrak g \longrightarrow \operatorname{End}_{K}(V)$，它使 $V$ 成为左 $\mathcal{U}\mathfrak g$-模. 由此得到 $\mathfrak g$ 的左模范畴与 $\mathcal{U}\mathfrak g$ 的左模范畴之间的范畴等价
\[\operatorname{Rep}(\mathfrak g) = \mathfrak g\text{-Mod} \xrightarrow{\ \sim\ } \mathcal{U}\mathfrak g\text{-Mod}.\]
特别地，$\mathfrak g$ 的不可约表示对应于 $\mathcal{U}\mathfrak g$ 的单模.
\begin{example*}
首先注意 $\mathcal{U}\mathfrak{gl}_{n}(K) \neq \mathrm{M}_{n}(K)$！！！例如取 $n = 2$：元素 $E = E_{12}$ 在 $\mathrm{M}_{2}(K)$ 中是 $2$ 阶幂零的；但在我们上面构造的表示 $V_{m}$ 中，它的像是 $m + 1$ 阶幂零的. 由于 $m$ 是任意的，这就蕴涵在 $\mathcal{U}\mathfrak{gl}_{2}$ 中它的像根本不是幂零的！事实上，$\mathcal{U}\mathfrak g$ 一般是无限维的.
\begin{enumerate}[label=\roman*)]
\item $\mathfrak g = K$. 可以验证 $\mathcal{U}\mathfrak g = K[T]$（其中 $\mathfrak g = K \xrightarrow{\ \sim\ } K \cdot T$）合乎要求.
\item $\mathfrak g$ 交换. 这时用来对 $T\mathfrak g$ 作商的理想由诸 $x \otimes y - y \otimes x$ 生成. 所得的代数是\newterm{ssl-symmetric-algebra}{对称代数}{symmetric algebra} $S\mathfrak g = \bigoplus_{n \in \mathbb{N}} S^{n}\mathfrak g$，它交换且分次，其构造适用于任何 $K$-向量空间. 函子 $V \mapsto SV$ 是``把交换代数映到其底层向量空间''的函子的左伴随. 有 $S^{n}\mathfrak g = (T^{n}\mathfrak g)_{\mathfrak{S}_{n}}$（对称群作用下的余不变量）. $\mathfrak g$ 的基的任一选取诱导同构 $S\mathfrak g \xrightarrow{\ \sim\ } K[T_{1}, \cdots, T_{\dim\mathfrak g}]$.
\item $\mathfrak g = \mathfrak{sl}_{2}(K)$. 可以给出如下的呈现：$\mathcal{U}\mathfrak{sl}_{2}$ 是由生成元 $e, f, h$ 及关系 $he - eh = 2e$、$hf - fh = -2f$ 与 $ef - fe = h$ 所定义的结合代数. 当 $K$ 的特征异于 $2$ 时，它包含一个值得注意的元素，即由
\[C := ef + fe + \frac{1}{2}h^{2}\]
给出的\newterm{ssl-casimir-operator}{Casimir 算子}{Casimir operator}.
\end{enumerate}
\end{example*}
\begin{exercise*}
仍设 $\mathfrak g$ 为 $\mathfrak{sl}_{2}$，且 $\operatorname{car}(K) \neq 2$.
\begin{enumerate}[label=\alph*)]
\item 证明 $C$ 与 $e$、$f$、$h$ 交换，从而属于 $\mathcal{U}\mathfrak g$ 的中心.
\item 若 $K$ 代数闭，推出 $C$ 在每个不可约表示上以纯量作用.
\item 证明在表示 $V_{m}$ 上，这个纯量为 $\lambda_{m} = \frac{1}{2}m(m + 2)$.
\item 证明 $C$ 在 $K[x, y]$ 上的作用（同~\ref{sec:ssl-2-1} 最后一个习题）由
\[\bigl( x\,\frac{\partial}{\partial x} + y\,\frac{\partial}{\partial y} \bigr) + \frac{1}{2} \bigl( x\,\frac{\partial}{\partial x} + y\,\frac{\partial}{\partial y} \bigr)^{2}\]
给出，并重新求出上述纯量.
\end{enumerate}
\end{exercise*}

\begin{numpara}[泛包络代数与微分算子]\label{np:ssl-2-2-2}
对光滑流形 $M$，记 $\operatorname{Diff}(\mathcal{C}^{\infty}(M))$ 为 $M$ 的（结合）\newterm{ssl-differential-operator}{微分算子}{differential operator}代数，亦即 $\operatorname{End}_{\mathbb{R}}(\mathcal{C}^{\infty}(M))$ 中由全体导子生成的 $\mathbb{R}$-子代数. 因此，一个微分算子局部上是形如 $f \cdot \frac{\partial^{r}}{\partial x_{1}^{r_{1}} \cdots \partial x_{n}^{r_{n}}}$（$f$ 为 $\mathcal{C}^{\infty}$ 函数）的算子的线性组合.
现在假设 $\mathfrak g = \operatorname{Lie}(G)$ 是某个 Lie 群 $G$ 的 Lie 代数（从而 $K = \mathbb{R}$），并考虑由左不变微分算子组成的子代数 $\operatorname{Diff}(\mathcal{C}^{\infty}(G))^{L}$. $\mathcal{U}\operatorname{Lie}(G)$ 的泛性质使我们能把映射 $X \in \operatorname{Lie}(G) \mapsto \xi_{X} \in \operatorname{Vec}(G)^{L}$ 扩张为代数态射
\[\mathcal{U}\operatorname{Lie}(G) \longrightarrow \operatorname{Diff}(\mathcal{C}^{\infty}(G))^{L}.\]
可以证明这个态射是一个同构.
\end{numpara}

\begin{numpara}[自由 Lie 代数]\label{np:ssl-2-2-3}
设 $I$ 是一个集合. 带有一组元素 $(X_{i})_{i \in I}$ 的 $K$-Lie 代数 $L$ 称为 $I$ 上的\newterm{ssl-free-lie-algebra}{自由 Lie 代数}{free Lie algebra}，如果对每个 $K$-Lie 代数 $\mathfrak g$，映射
\[\text{($*$)}\qquad \operatorname{Hom}_{K\text{-Lie}}(L, \mathfrak g) \longrightarrow \mathfrak g^{I} = \operatorname{Hom}_{\mathrm{Ens}}(I, \mathfrak g),\qquad f \longmapsto (f(X_{i}))_{i \in I}\]
是双射. 这等价于说 $(L, (X_{i})_{i \in I})$ 是``$K$-Lie 代数与以 $I$ 为指标的一族元素''这类二元组所成范畴中的始对象. 鉴于条件中的泛性，$I$ 上的自由 Lie 代数在唯一同构意义下唯一.
现在来证明它存在. 记 $V := \bigoplus_{i \in I} K X_{i}$ 为以 $(X_{i})_{i \in I}$ 为基的向量空间. 张量代数 $TV$ 显然是 $I$ 上的自由结合代数（定义与上面类似）. 特别地，对每个 $K$-Lie 代数 $\mathfrak g$，映射
\[\operatorname{Hom}_{K\text{-alg}}(TV, \mathcal{U}\mathfrak g) \longrightarrow \mathcal{U}\mathfrak g^{I} = \operatorname{Hom}_{\mathrm{Ens}}(I, \mathcal{U}\mathfrak g),\qquad f \longmapsto (f(X_{i}))_{i \in I}\]
是双射. 把 $L$ 定义为 $TV$ 中由（$1$ 次的）元素 $X_{i}$（$i \in I$）生成的 Lie 子代数. 映射 $(*)$ 这时显然是单射（因为诸 $X_{i}$ 生成 $L$）. 为证满射，暂时承认 $\mathfrak g \longrightarrow \mathcal{U}\mathfrak g$ 是单射（见下一段），并把 $\mathfrak g$ 等同于 $\mathcal{U}\mathfrak g$ 的一个 Lie 子代数. 任取 $\mathfrak g \subset \mathcal{U}\mathfrak g$ 中的一族元素 $(Y_{i})_{i \in I}$，设由上述双射给出的代数态射为 $f : TV \longrightarrow \mathcal{U}\mathfrak g$. 那么 $f(V) \subset \mathfrak g$，因此 $f(L) \subset \mathfrak g$ 且 $f(X_{i}) = Y_{i}$，如愿.
\end{numpara}
\begin{exercise*}
证明 $I$ 上的自由 Lie 代数的\newterm{ssl-enveloping-algebra}{泛包络代数}{universal enveloping algebra}是 $I$ 上的自由结合代数.
\end{exercise*}

\begin{numpara}[Poincaré--Birkhoff--Witt 定理]\label{np:ssl-2-2-4}
回到满射 $\pi : T\mathfrak g \twoheadrightarrow \mathcal{U}\mathfrak g$. 若给 $T\mathfrak g$ 赋予滤过 $T\mathfrak g^{\leqslant m} := \bigoplus_{k \leqslant m} T^{k}\mathfrak g$，并给 $\mathcal{U}\mathfrak g$ 赋予像滤过，则 $\pi$ 是滤过结合代数的态射. 过渡到相伴分次，便得到一个分次代数的、仍是满射的态射
\[\operatorname{gr}\pi : T\mathfrak g \xrightarrow{\ \sim\ } \operatorname{gr}(T\mathfrak g) := \bigoplus_{k \in \mathbb{N}} T\mathfrak g^{\leqslant k}/T\mathfrak g^{\leqslant k-1} \longrightarrow \operatorname{gr}(\mathcal{U}\mathfrak g) := \bigoplus_{k \in \mathbb{N}} \mathcal{U}\mathfrak g^{\leqslant k}/\mathcal{U}\mathfrak g^{\leqslant k-1}.\]
由于对 $x, y \in \mathfrak g$ 有 $\pi(x \otimes y - y \otimes x) = \pi([x, y])$，所以也有 $\operatorname{gr}\pi(x \otimes y) = \operatorname{gr}\pi(y \otimes x)$. 由此可知，$\operatorname{gr}\pi$ 经过对称代数（见上文的例 ii））分解为分次代数的一个满态射 $S\mathfrak g \longrightarrow \operatorname{gr}(\mathcal{U}\mathfrak g)$；特别地，$\operatorname{gr}(\mathcal{U}\mathfrak g)$ 是交换的.
\begin{theorem*}
（PBW，第一种表述）上述态射 $S\mathfrak g \xrightarrow{\ \operatorname{gr}\pi\ } \operatorname{gr}(\mathcal{U}\mathfrak g)$ 是单射. 因而它是分次 $K$-代数的同构.
\end{theorem*}
出于时间原因，我们不在此证明这个定理（例如可参看 Carter 的书或 Humphreys 的书）. 下面给出它的一个有意思的推论作为例子.
\begin{corollary*}
\begin{enumerate}[label=\roman*)]
\item 映射 $\iota : \mathfrak g \hookrightarrow \mathcal{U}\mathfrak g$ 是单射.
\item 对 $\mathfrak g$ 的任何 Lie 子代数 $\mathfrak h$，相应的态射 $\mathcal{U}\mathfrak h \longrightarrow \mathcal{U}\mathfrak g$ 是单射.
\end{enumerate}
\end{corollary*}
注意，当 $\mathfrak g$ 有限维时，i) 也可以由 Ado 定理推出，该定理断言 $\mathfrak g$ 可嵌入某个 $\mathfrak{gl}(V)$，从而嵌入 $\operatorname{End}_{K}(V)$.
下面给出上述定理在 $\mathfrak g$ 有限维的情形下更具体但精确度稍逊的一种表述.
\begin{theorem*}
（PBW，第二种表述）设 $x_{1}, \cdots, x_{n}$ 是 $\mathfrak g$ 在 $K$ 上的一组基. 则形如 $x_{i_{1}}x_{i_{2}} \cdots x_{i_{k}}$（其中 $i_{1} \leqslant i_{2} \leqslant \cdots \leqslant i_{k}$）的单项式构成 $\mathcal{U}\mathfrak g$ 在 $K$ 上的一组基.
\end{theorem*}
在无限维的情形，选取 $\mathfrak g$ 的一个全序基，可以得到类似的陈述.
\end{numpara}

\begin{numpara}[双代数结构]\label{np:ssl-2-2-5}
最后作如下注记，它表明一个泛包络代数携带的结构比结合代数的更多. 考虑把 $x$ 映到 $\iota(x) \otimes 1 + 1 \otimes \iota(x)$ 的映射 $\mathfrak g \longrightarrow \mathcal{U}\mathfrak g \otimes_{K} \mathcal{U}\mathfrak g$. 可以证明它是 Lie 代数态射（习题）. 因此它诱导代数态射 $\Delta : \mathcal{U}\mathfrak g \longrightarrow \mathcal{U}\mathfrak g \otimes_{K} \mathcal{U}\mathfrak g$. 进而可以证明恒等式 $(\Delta \otimes \operatorname{id}) \circ \Delta = (\operatorname{id} \otimes \Delta) \circ \Delta$（余结合性公理）. 我们说 $\Delta$ 是一个\newterm{ssl-coproduct}{余积}{coproduct}，而 $\mathcal{U}\mathfrak g$ 是一个\newterm{ssl-bialgebra}{双代数}{bialgebra}（甚至是一种``Hopf 代数''）. 我们在这里不再对此多作讨论，但这一观点在\newterm{ssl-quantum-group}{量子群}{quantum group}理论中是根本性的. 后者并不是真正意义上的群，而是经典泛包络代数的一种``形变''.
\end{numpara}
